\documentclass{article}
% if you need to pass options to natbib, use, e.g.:
%     \PassOptionsToPackage{numbers, compress}{natbib}
% before loading neurips_2022


% ready for submission
\usepackage[preprint]{neurips_2026}


% to compile a preprint version, e.g., for submission to arXiv, add add the
% [preprint] option:
%     \usepackage[preprint]{neurips_2022}


% to compile a camera-ready version, add the [final] option, e.g.:
%     \usepackage[final]{neurips_2022}


% to avoid loading the natbib package, add option nonatbib:
%    \usepackage[nonatbib]{neurips_2022}


\usepackage[utf8]{inputenc} % allow utf-8 input
\usepackage[T1]{fontenc}    % use 8-bit T1 fonts
\usepackage{hyperref}       % hyperlinks
\usepackage{url}            % simple URL typesetting
\usepackage{booktabs}       % professional-quality tables
\usepackage{amsfonts}       % blackboard math symbols
\usepackage{nicefrac}       % compact symbols for 1/2, etc.
\usepackage{microtype}      % microtypography
\usepackage{xcolor}         % colors
\newcommand{\new}[1]{{\color{red} #1}}

\title{Refining Video Object Segmentation with Test-time Gradient Correction}

\author{%
  Suyu Chen \\
  \texttt{suc033@ucsd.edu}
  \And
  Yi Sun \\
  \texttt{yis094@ucsd.edu}
  \And
  Yinru Chen\\
  \texttt{yic173@ucsd.edu}
  \AND
  \textbf{Team \#71}
}

\begin{document}

\maketitle

\begin{abstract}
Provide a concise summary of your project. Your abstract should briefly describe the problem, why it is important, your main approach, key results, and the main conclusion. Note your track here. 
\end{abstract}

\section{Introduction}

Clearly introduce the problem your project addresses and explain why it is important. Provide enough background for readers to understand the motivation, then summarize your project goal, approach, and main contributions.

\section{Related Work}

Discuss the most relevant prior work related to your project. Rather than listing papers one by one, try to organize them by themes, explain their limitations, and clearly describe how your project is related to or different from existing work.

\section{Method / Approach}

Describe your proposed method or approach in detail. Include the problem formulation, model architecture, algorithms, equations, or technical design when appropriate, and explain why your approach is suitable for the problem.

\section{Results \& Analysis}

Present your experimental results clearly using tables, figures, and/or qualitative examples. Include sufficient implementation details, such as datasets, preprocessing, train/test splits, baselines, evaluation metrics, and hyperparameters, so that your results are reproducible. Analyze the results carefully and explain what they demonstrate about your method.

\section{Conclusion}

Summarize the main findings of your project and reflect on what you learned. Discuss the limitations of your current work and suggest meaningful future directions.

\section{Code \& GitHub}

Provide a link to your GitHub repository or code submission. Your repository should be well organized and include a clear README with instructions for running the code and reproducing the main results.

\section{Team Member Contribution}

Clearly describe each team member's contribution to the project. Please be specific about each person's role in literature review, method development, implementation, experiments, analysis, writing, and presentation.

\section*{Bonus: Course Evaluation Submission}

We hope you enjoyed the course, and we would really appreciate you taking a moment to complete the teaching evaluation: 
\url{https://academicaffairs.ucsd.edu/Modules/Evals/}

The deadline is Saturday, June 6 at 8:00 AM. To receive the 5-point bonus, please include a screenshot of your evaluation submission in the Appendix after the “Team Member Contribution” section. You do not need to show the contents of your evaluation. The bonus will be added to your Final Project Report score, but the total report score will be capped at 80 points.

\end{document}
